\documentclass[11pt]{article}
\usepackage[margin=1in]{geometry}
\usepackage{amsmath,amssymb}
\usepackage{booktabs}
\usepackage{graphicx}
\usepackage[hidelinks]{hyperref}
\usepackage{xcolor}

\newcommand{\gw}{\textsc{gradwave}}
\newcommand{\qe}{\textsc{Quantum ESPRESSO}}
\newcommand{\meV}{\,\mathrm{meV}}
\newcommand{\ueV}{\,\mu\mathrm{eV}}

\title{\gw: a differentiable plane-wave DFT code in PyTorch\\
\large Capability portfolio: performance and automatic differentiation}
\author{Will Laderer}
\date{July 2026}

\begin{document}
\maketitle

\begin{abstract}
\gw{} is a plane-wave pseudopotential DFT code built on PyTorch, designed so
that every physical quantity is a differentiable function of its inputs.
It reproduces \qe{} 7.5 at mature-code precision (equation-of-state
$\Delta$-factor $\leq 0.17\meV$/atom over six materials; the full
QE-comparison test suite passes on both CPU and a consumer GPU), while
exposing capabilities that conventional codes cannot: exact gradients of the
total energy with respect to exchange--correlation parameters, Hubbard $U$,
and atomic positions, obtained from the computation graph rather than
hand-derived matrix elements. This document collects a wall-clock comparison
against \qe{} and four automatic-differentiation showcases: learnable
exchange--correlation functionals, the Hubbard $U$ as a differentiable and
self-determinable parameter, phonons from autograd forces, and noncollinear
magnetism with the exchange field generated by autograd.
\end{abstract}

\section{Overview and validation summary}

\gw{} implements norm-conserving pseudopotential DFT (UPF~v2, ONCV/PseudoDojo)
with LDA/PBE and collinear/noncollinear spin including spin--orbit coupling,
in float64/complex128 throughout. The code is organized in three layers:
a pure differentiable tensor core (energies, Hamiltonian applies), a
setup/solver layer (Davidson, SCF mixing) that runs under
\texttt{torch.no\_grad()}, and a user layer (YAML/CLI, ASE calculator).
Autograd never traces the SCF loop; derivatives are obtained at the converged
fixed point via stationarity (Hellmann--Feynman), implicit differentiation
(Sternheimer response with the Hxc kernel as an autograd Hessian-vector
product), or explicit graphs over the position-dependent terms.

\begin{table}[h]\centering
\caption{Headline agreement with \qe{}~7.5 at identical pseudopotentials,
cutoffs, $k$-meshes, and FFT grids.}
\begin{tabular}{llll}
\toprule
benchmark & quantity & \gw{} vs \qe{} & scale \\
\midrule
SCF energies (Si, Al, Cu, Fe, GaAs, \ldots) & $E$ or $F$ & $\leq 2\meV$/atom (most $\leq 10^{-3}$) & --- \\
$\Delta$-factor EOS (6 materials, $\pm 6\%$) & $\Delta$ & $\leq 0.17\meV$/atom & mature codes: 0.3--1 \\
lattice constants from EOS & $a_0$ & $\leq 4\times10^{-4}$\,\AA & --- \\
forces (displaced Si, Al) & $F$ & $\leq 5\meV$/\AA & --- \\
NVE drift (Si$_8$, 0.5\,ps, autograd forces) & $E_{\rm tot}$ secular & $\sim 1\ueV$/atom & exact-force floor \\
dielectric constant (Si) & $\varepsilon_\infty$ & 0.002\,\% & vs \texttt{ph.x} DFPT \\
Born charges (MgO) & $Z^*$ & 0.07--0.10\,\% & vs \texttt{ph.x} DFPT \\
band structures (Si) & $\epsilon_{n\mathbf{k}}$ & $\leq 10\meV$ & --- \\
SOC splitting (GaAs $\Delta_0$) & $0.3360$ vs $0.3360$\,eV & $0.03\meV$ & spinor + $j$-projectors \\
NiO DFT+$U$ (AFM-II) & $E_U$ & $1.6\meV$ & occupations $10^{-3}$ \\
\bottomrule
\end{tabular}
\end{table}

The complete integration suite (46 tests, including every slow QE comparison)
passes unmodified on CUDA; a rattled-Si$_8$ NVE trajectory driven by the ASE
calculator conserves total energy to a secular drift of
$\sim 1\ueV$/atom over 0.5\,ps --- direct evidence that the autograd
Hellmann--Feynman force is the \emph{exact} gradient of the SCF energy
surface.

\section{Performance versus Quantum ESPRESSO}
\label{sec:speed}

Wall-clock comparison on a single machine (Intel laptop, 22 hardware threads,
NVIDIA RTX~3050~6\,GB). \qe{}~7.5 runs with \texttt{mpirun -np 8} (its
efficient parallelization sweet spot for these small cells); \gw{} runs on
22 CPU threads or on the GPU. Both codes use identical plane-wave cutoffs,
$k$-meshes (symmetry reduction on in both), FFT grids (pinned), band counts,
functionals, and comparable convergence thresholds
(\texttt{conv\_thr}~$=10^{-10}$\,Ry vs \texttt{etol}~$=10^{-9}$\,eV).
Times are full SCF wall times excluding setup.

\begin{table}[h]\centering
\caption{SCF wall time (seconds); identical physics settings, energies agree
to $\lesssim 1\meV$/atom. Iteration counts in parentheses.}
\begin{tabular}{lrrrrrr}
\toprule
system & $N_{\rm at}$ & ecut (Ry) & $k$-mesh & \qe{} (8 MPI) & \gw{} CPU-22 & \gw{} GPU \\
\midrule
Si (dia.)   & 2 & 30 & $4^3$ & 1.5 \,(6) & 3.5 \,(10) & 1.5 \,(10) \\
C (dia.)    & 2 & 40 & $4^3$ & 1.5 \,(6) & 1.8 \,(11) & 0.9 \,(11) \\
GaAs        & 2 & 40 & $4^3$ & 1.8 \,(9) & 7.9 \,(14) & 2.8 \,(14) \\
Al (fcc)    & 1 & 40 & $8^3$ & 1.8 \,(6) & 5.1 \,(11) & 4.2 \,(11) \\
Cu (fcc)    & 1 & 45 & $8^3$ & 2.2 \,(9) & 10.0 \,(13) & 5.1 \,(13) \\
MgO         & 2 & 50 & $4^3$ & 1.6 \,(9) & 3.6 \,(13) & 1.6 \,(13) \\
Si$_8$      & 8 & 30 & $2^3$ & 1.9 \,(8) & 5.8 \,(13) & 5.4 \,(13) \\
\bottomrule
\end{tabular}
\end{table}

The honest summary: on cells this small, \qe{}'s two decades of MPI
engineering show --- it is on par with or faster than \gw{}'s GPU path and
2--5$\times$ faster than the 22-thread CPU path. Two structural factors
account for most of the gap: \qe{} converges in 6--9 SCF iterations to
\gw's 10--14 (mixer maturity, not tensor throughput --- per-iteration GPU
cost is at or below \qe's), and at these sizes ($\leq 3000$ plane waves)
kernel-launch and Python overhead are proportionally largest. The GPU
advantage grows with system size (e.g.\ a 78-electron Bi$_2$Se$_3$ SOC
calculation on an $80^3$ grid runs its SCF in minutes on the 6\,GB card);
the relevant point for this portfolio is that the fully differentiable
formulation costs little enough that the same code base is usable for
production ground states \emph{and} everything in the following sections.

\section{Differentiable exchange--correlation functionals}
\label{sec:xc}

The XC functional is an \texttt{nn.Module}; $v_{\rm xc} = \delta E_{\rm
xc}/\delta\rho$ is generated by autograd (including GGA gradient terms), so
any twice-differentiable functional --- including one with trainable
parameters --- is a drop-in. At SCF convergence the energy is variational in
the density, so the training gradient is free:
\begin{equation}
\frac{dE}{d\theta} = \left.\frac{\partial E_{\rm xc}[\rho;\theta]}
{\partial\theta}\right|_{\rho^\star} ,
\end{equation}
one backward pass, no response solve, no finite differences.
Table~\ref{tab:xc} verifies this against central finite differences of full
SCF re-runs for a PBE-form exchange with learnable $(\kappa,\mu)$
(initialization at PBE values reproduces PBE to $10^{-12}$).

\begin{table}[h]\centering
\caption{$dE/d\theta$ for learnable exchange parameters (Si, off-PBE point
$\kappa=0.65$, $\mu=0.18$): one autograd backward vs.\ SCF finite
differences.}
\label{tab:xc}
\begin{tabular}{lrrr}
\toprule
parameter & autodiff & finite difference & $|$diff$|$ \\
\midrule
$\theta_\kappa$ & $-0.1911153641$ & $-0.1911153645$ & $3.3\times10^{-10}$ \\
$\theta_\mu$    & $-1.6478174399$ & $-1.6478174406$ & $6.7\times10^{-10}$ \\
\bottomrule
\end{tabular}
\end{table}

As an end-to-end demonstration, starting from an off-PBE point
$(\kappa,\mu) = (0.55, 0.30)$ and fitting to PBE total energies at two
volumes with Adam:
\begin{center}
loss $6.5\times10^{-1} \rightarrow 3.8\times10^{-4}$ in five steps;\quad
$\mu \rightarrow 0.2222$ vs PBE's $0.2195$ after ten steps.
\end{center}
The energy data pin down $\mu$ (in silicon's reduced-gradient range,
$F \approx 1 + \mu s^2$ to leading order) while $\kappa$ --- an $s^4$
effect --- is left nearly unconstrained and wanders a flat valley: the
optimization itself performs a physically meaningful identifiability
analysis of the functional form. Density-based losses use the implicit SCF backward
(self-consistent adjoint/Sternheimer, validated to $2\times10^{-4}$ against
finite differences).

\section{The Hubbard $U$ as a differentiable quantity}
\label{sec:hubbardu}

\gw{} implements rotationally invariant (Dudarev) DFT+$U$ with QE's
\texttt{atomic} projector convention, validated on AFM-II NiO against \qe{}
($E_U$: 2.9956 vs 2.9942\,eV; Ni-$d$ occupations 4.909/3.647 vs 4.910/3.648;
$E(U{=}5) - E(U{=}0)$: 3.3245 vs 3.3216\,eV). On top of this, three
derivative capabilities:

\paragraph{Exact $dE/dU$.} By stationarity, $dE_{\rm tot}/dU = \sum_{I\sigma}
\tfrac12\,{\rm Tr}\!\left[n^{I\sigma}(1-n^{I\sigma})\right]$ at the converged
point --- the gradient a $U$-learning loop would backpropagate. NiO at
$U=5$\,eV: analytic $0.616154$ vs SCF finite differences $0.616153$
($|{\rm diff}| = 7.4\times10^{-7}$).

\paragraph{$+U$ forces by autograd.} $E_U$ depends on positions only through
the projector phases $e^{-i(\mathbf{k}+\mathbf{G})\cdot\tau}$; autograd
through the phase factors gives $-dE_U/d\tau$ directly. On symmetry-broken
NiO the analytic force matches frozen-orbital finite differences to
$4.9\times10^{-8}$\,eV/\AA{} and vanishes identically on atoms without a
correlated manifold.

\paragraph{The code computes its own $U$.} The Cococcioni--de Gironcoli
linear-response $U = (\chi_0^{-1} - \chi^{-1})_{II}$ is implemented two ways.
The finite-difference route ($\pm\alpha$ probe SCF re-runs) gives
$U = 6.4493$\,eV for NiO. The \emph{analytic} route obtains both response
matrices from a single ground state: a batched conduction-projected
Sternheimer solve with the projector probe on the right-hand side, screened
self-consistently through the spin Hxc kernel evaluated as an autograd
Hessian-vector product of $E_{\rm Hxc}$ --- DFPT without hand-coding
$f_{\rm xc}^{\sigma\sigma'}$, valid for any twice-differentiable spin
functional including learnable ones. Anderson acceleration is essential: the
antisymmetric magnetization mode across the two Ni sites gives
$K\chi_0$ an eigenvalue $\approx -6$ (PBE NiO is near a spin instability),
which diverges under plain damped iteration; Anderson converges in 11
iterations.

\begin{table}[h]\centering
\caption{Linear-response Hubbard $U$ for NiO (Ni $3d$, PBE, PseudoDojo, 40\,Ry,
$2^3$ $k$-mesh).}
\begin{tabular}{lr}
\toprule
method & $U$ (eV) \\
\midrule
\gw{}, finite-difference response ($\pm\alpha$ SCFs) & 6.4493 \\
\gw{}, analytic Sternheimer + autograd-HVP kernel & 6.4496 \\
\qe{} \texttt{hp.x} (DFPT, $n_q = 1$) & 6.4308 \\
\bottomrule
\end{tabular}
\end{table}

The two \gw{} routes agree to $0.3\meV$ (and column-by-column to
$2\times10^{-5}$ in $\chi$, $\chi_0$); both agree with \qe's independent
DFPT implementation to $0.3\%$.

\section{Phonons from autograd forces}
\label{sec:phonons}

Forces are exact autograd gradients, so force-constant matrices from finite
displacements inherit machine-grade consistency: the Hessian entry from
displaced \emph{forces} matches the second difference of \emph{energies} to
$<0.5\%$, acoustic modes vanish under the sum rule, and the NVE test of
Sec.~1 bounds any systematic force error at the $\ueV$ scale.
At settings identical to a \qe{} \texttt{ph.x} DFPT calculation
(PBE, 40\,Ry, $4^3$ mesh, pinned $32^3$ grid):

\begin{table}[h]\centering
\caption{Si $\Gamma$ phonons (cm$^{-1}$): finite differences of \gw{}
autograd forces vs \qe{} \texttt{ph.x} DFPT.}
\begin{tabular}{lrr}
\toprule
mode & \gw{} (FD of autograd forces) & \qe{} \texttt{ph.x} (DFPT) \\
\midrule
acoustic $T_{1u}$ ($\times 3$) & $0.000$ (ASR) & $6.03$ (raw, no ASR) \\
optical $T_{2g}$ ($\times 3$)  & $522.55$ & $523.35$ \\
\bottomrule
\end{tabular}
\end{table}

The optical frequency agrees with the independent DFPT result to
$0.8$\,cm$^{-1}$ ($0.15\%$), the triplet is degenerate to the printed
precision, and the acoustic modes vanish exactly under the sum rule ---
all from displacements of $5\times10^{-3}$\,\AA{} evaluated with autograd
forces, with no analytic force implementation anywhere in the code.

For polar materials the ingredients for LO--TO splitting are available from
the same response stack: $\varepsilon_\infty$ and Born charges
(Sec.~1) are computed by E-field DFPT in which the Born charge is literally
a mixed derivative $\partial^2 E/\partial\mathcal{E}\,\partial\tau$ ---
one autograd backward through the position-differentiable pseudopotential
per field direction yields $Z^*$ for all atoms simultaneously.

\section{Noncollinear magnetism with autograd exchange fields}
\label{sec:ncl}

The noncollinear (spinor) SCF uses locally collinear XC in which a
\emph{single} autograd call on $E_{\rm xc}[\rho, |\mathbf{m}|]$ yields both
$v_{\rm xc}$ and the exchange field $\mathbf{B}_{\rm xc} \parallel \mathbf{m}$
--- no hand-derived functional derivatives of the spin rotation. Validation
ladder on bcc Fe:

\begin{itemize}
\item collinear limit: spinor SCF with $\mathbf{m} \parallel \hat z$
  reproduces the two-channel calculation to $0.0000\meV$;
\item rotational invariance: total energies for $\mathbf{m}$ along
  $\hat x$, $\hat z$, and $45^\circ$ agree to $0.19\ueV$ (no SOC ---
  the exact symmetry emerges numerically from the autograd field);
\item moment directions are preserved self-consistently (no drift to a
  collinear attractor).
\end{itemize}

With fully relativistic $j$-resolved projectors (spin--orbit coupling built
from complex spherical harmonics and Clebsch--Gordan spin structure on the
doubled plane-wave axis):

\begin{itemize}
\item GaAs: free energy matches \qe{} \texttt{lspinorb} to $0.0007\meV$/atom;
  split-off gap $\Delta_0 = 0.3360$ vs 0.3360\,eV ($0.03\meV$);
\item Bi$_2$Se$_3$: the topological band inversion is reproduced --- the
  $\Gamma$ gap-edge parities swap between scalar-relativistic and SOC
  calculations, with parity eigenvalues $\pm 1.000$ and exact Kramers
  degeneracy (78 electrons, $80^3$ grid, NLCC, on the RTX 3050);
\item magnetic energy mapping: bcc Fe Heisenberg couplings
  $J_1 = +57.5\meV$, $J_2 = -39.6\meV$ (mean-field $T_c \approx 860$\,K)
  from signed-shell energy differences; Cr correctly orders AFM
  ($-27.7\meV$/atom vs nonmagnetic, matching \qe{} to $0.02\meV$/atom).
\end{itemize}

The same differentiable structure is what makes magnetic \emph{response}
quantities tractable next: the spin Hxc kernel used for the Hubbard-$U$
response (Sec.~\ref{sec:hubbardu}) is the $2\times2$ spin block of exactly
this autograd machinery, and moment-direction torques $dE/d\hat{\mathbf{m}}$
for magnetocrystalline anisotropy follow the same pattern with SOC.

\section{Closing remarks}

Every result above ran on commodity hardware (a laptop CPU and a 6\,GB
consumer GPU) against reference calculations from \qe{}~7.5 with identical
numerical settings. The differentiable formulation did not cost ground-state
precision (Sec.~1) or speed (Sec.~\ref{sec:speed}); it added a family of
exact derivatives --- $dE/d\theta_{\rm xc}$, $dE/dU$, $dE/d\tau$,
$\partial^2 E/\partial\mathcal{E}\partial\tau$, self-consistent response via
autograd Hxc kernels --- that conventional codes obtain only through
hand-derived, functional-specific implementations, if at all.

\end{document}
